% ==================== 五、问题三 ====================
\section{问题三：预报融合与场景树多时点调整}

\subsection{问题分析}

问题三在问题二基础上引入两项新要素：

\begin{enumerate}[itemsep=3pt,topsep=3pt]
  \item \textbf{多时点预报}：每天 $0$、$6$、$12$、$18$ 时各发布一次未来 24 h 的光伏整点预报，
        后续时刻的预报可用于\textbf{调整}购电策略；
  \item \textbf{调整计费}：调整不是免费的——
        计划量高于调整量的部分按 $50\%$ 违约电价、调整量高于计划量的部分按 $1.5$ 倍电价。
\end{enumerate}

因此模型需要回答两个层次的问题：
\textbf{(i) 如何在 0 点制定“预留调整余地”的计划：}
\textbf{(ii) 如何在后续发布时刻利用新信息最优地修正计划：}

\subsection{计费规则的统一化建模}
\label{sec:q3-fee}

设第 $t$ 段的 0 点原计划量为 $g_t^0$，最终调整后的购电量为 $q_t$。定义调增量与调减量

\begin{equation}
u_t=(q_t-g_t^0)_+,\qquad v_t=(g_t^0-q_t)_+,
\end{equation}

则 $q_t=g_t^0+u_t-v_t$。题面规则可统一写为

\begin{equation}
J=\sum_t\Bigl[\underbrace{p_tg_t^0}_{\text{原计划费}}
+\underbrace{1.5p_tu_t}_{\text{调增，1.5 倍}}
-\underbrace{p_tv_t}_{\text{退还原价}}
+\underbrace{0.5p_tv_t}_{\text{违约费，50\%}}
+\underbrace{5p_te_t}_{\text{紧急购电}}\Bigr].
\label{eq:q3-fee}
\end{equation}

即\textbf{净调整费为 $1.5p_tu_t-0.5p_tv_t$，它可以为负}。
合并同类项得另一便于独立核验的等价式：

\begin{equation}
J=\sum_t\Bigl[p_t\min(g_t^0,q_t)+0.5p_tv_t+1.5p_tu_t+5p_te_t\Bigr].
\label{eq:q3-fee2}
\end{equation}

\subsection{在线预报融合}

\subsubsection{融合权重}
\label{sec:q3-fusion}

设历史预测为 $H$、发布预报为 $F$、实际值为 $V$。只用最近 56 天内\textbf{已完全实现}的预报行，
以最小二乘拟合权重：

\begin{equation}
\alpha=\operatorname{clip}_{[0,1]}\frac{\sum(F-H)(V-H)}{\sum(F-H)^2},
\qquad \text{融合值}=\alpha F+(1-\alpha)H.
\label{eq:q3-alpha}
\end{equation}

\textbf{式 \eqref{eq:q3-alpha} 的推导}：以融合值与实际值的平方误差为目标
$\min_{\alpha}\sum\bigl[V-(\alpha F+(1-\alpha)H)\bigr]^2$。
记 $z=F-H$，则中括号内为 $(V-H)-\alpha z$，目标化为
$f(\alpha)=\sum\bigl[(V-H)-\alpha z\bigr]^2$。
对 $\alpha$ 求导并置零：
\begin{equation}
\frac{\mathrm{d}f}{\mathrm{d}\alpha}
=-2\sum z\bigl[(V-H)-\alpha z\bigr]=0
\ \Longrightarrow\ \alpha=\frac{\sum z\,(V-H)}{\sum z^2}.\label{eq:q3-alpha-deriv}
\end{equation}
又因二阶导 $\mathrm{d}^2f/\mathrm{d}\alpha^2=2\sum z^2>0$，
该驻点即全局极小点；再叠加 $[0,1]$ 截断以避免外推。
几何上，$\alpha$ 是 $V-H$ 在 $F-H$ 上的投影系数。

\textbf{收缩偏差校正}：进一步叠加历史平均残差的收缩估计
\begin{equation}
\text{偏差}=\overline{(V-\text{融合})}\cdot\frac{n}{n+14},
\end{equation}
其中 $n$ 为有效样本数。系数 $n/(n+14)$ 用于抑制小样本下偏差估计的波动：
样本充足时接近 $1$，样本稀少时被压缩向 $0$。

\subsubsection{发布预报的时间对齐}

附件 3 的“预报 $k$ 小时”是\textbf{发布时刻后第 $k$ 个整点}的\textbf{瞬时}功率，
而非第 $k$ 个小时区间的均值。故采用如下重建方式：

\begin{enumerate}[itemsep=2pt,topsep=3pt]
  \item 以 24 个整点为锚点，在 10 分钟网格上\textbf{线性插值}；
  \item 发布瞬间的端点使用\textbf{已观测}的实际光伏功率（0 点用前一日末观测值）。
\end{enumerate}

\noindent\textbf{为何不能把整点值当作“小时均值”：} 设真实光伏功率为正弦形的日间过程，
其峰谷差为 $A$、周期 $T_0=24$ h。若把第 $k$ 个整点的值 $v(k)$ 当作区间
$[k,k+1)$ 的均值作分段常数填充，相当于用左端点值代替中点值，
引入的按时段均值计的系统偏差为
\begin{equation}
\overline{\varepsilon}=\frac{1}{T_0}\int_0^{T_0}\bigl[v(k)-v(t)\bigr]\,\mathrm{d}t\ \sim\ \frac{1}{2}\cdot\frac{\mathrm{d}v}{\mathrm{d}t}\cdot 1\ \text{h},
\label{eq:q3-phase}
\end{equation}
即在梯度最大处（日出、日落）偏差最大，量级为\textbf{半小时}的光伏变化量，
故称 “$1/2$ 小时量级的系统相位误差”。线性插值则以相邻两整点为端点，
把该项偏差从 $O(\Delta)$ 降到 $O(\Delta^2)$，且在每个整点上无偏。

\noindent\textbf{实证核对：} 以附件 2 的实际光伏为基准，用两种方式重建同一预报行的
144 段并逐段比较（除去最后一天以避开跨年缺口），结果见表 \ref{tab:q3-align}。

\begin{table}[H]
\centering
\caption{两种整点重建方式的全年逐段 MAE（kW，365 天平均）}
\label{tab:q3-align}
\begin{tabular}{@{}l r r r@{}}
\toprule
发布时间 & 分段常数填充 & \textbf{线性插值（本文）} & 降幅 \\
\midrule
$0$ 点  & 440.64 & \textbf{225.99} & $-48.7\%$ \\
$6$ 点  & 405.15 & \textbf{167.39} & $-58.7\%$ \\
$12$ 点 & 453.61 & \textbf{238.10} & $-47.5\%$ \\
$18$ 点 & 488.43 & \textbf{292.97} & $-40.0\%$ \\
\bottomrule
\end{tabular}
\end{table}

\noindent 线性插值在四个发布时间上把逐段 MAE 降低 $40.0\%$--$58.7\%$，
与式 \eqref{eq:q3-phase} 的机理一致：分段常数填充在日出、日落的高梯度处偏差最大，
而线性插值恰在这些位置把误差从 $O(\Delta)$ 降到 $O(\Delta^2)$。
这说明\textbf{时间对齐方式本身是预报误差的一个可测量来源}，而非实现细节。
需要说明的是，表 \ref{tab:q3-align} 度量的是“给定整点预报下的重建误差”，
而图 \ref{fig:q3-fusion} 展示的是融合后的综合预报精度，两者口径不同、不可直接相减。

\subsection{经验场景树}

\subsubsection{信息限制：非前视约束（本问的建模前提）}
\label{sec:q3-nonanticip}

问题二已要求“预测只能用历史数据”，
本问把它\textbf{进一步强化}到“场景树的节点划分也只能依据届时已发布的信息”，
原因是本问新增了多个决策时点，若不限制就可能在 6 点“提前知道” 12 点才会发布的信息。
具体为三条：
\begin{itemize}[itemsep=2pt,topsep=3pt]
  \item 训练样本仅取\textbf{不晚于 $d-2$ 天发布}的预报行，
        以确保跨午夜的未来 24 h 已经实现；
  \item 每个发布时刻\textbf{分别}训练，不得利用当天后续发布的预报来估计当前权重；
  \item 样本少于 7 条时退化为直接使用发布预报。
\end{itemize}

\noindent 上述约束在模型中的实现方式是“节点内共享变量”，
见 \S\ref{sec:q3-tree}；它比额外添加约束更稳健，因为物理上不可能针对已知场景分别决策。

\subsubsection{场景构造}

在第 $a$ 次发布时刻，场景净负载为

\begin{equation}
N_{s,t}=\hat N_t^{(a)}+\Bigl(N_{r_s,t}-\hat N_{r_s,t}^{(a)}\Bigr),
\label{eq:q3-scen}
\end{equation}

即用历史日 $r_s$ 在\textbf{同一发布时刻}的残差去修正当前的发布预报。

\subsubsection{信息分支（场景树）}
\label{sec:q3-tree}

仅用点预测场景会丢失一个重要信息：\textbf{不同历史日不仅能告诉我们“缺口可能多大”，
还能告诉我们“预报会如何变化”}。为此把场景组织为树：

\begin{enumerate}[itemsep=3pt,topsep=3pt]
  \item 历史日 $r_s$ 除净负载残差外，还附带它在\textbf{后续发布时刻}的预报修订向量；
  \item 把这些修订作为“未来可能观测到的信息”，用\textbf{二叉聚类}（按方差最大的维度二分）
        组织成树；
  \item \textbf{关键的非前视要求}：每个节点只能依据\textbf{届时新发布的预报修订}划分其子节点，
        不得使用尚未实现的未来负载或实际光伏；
  \item 节点概率 $\pi_n=$ 落入该节点的历史日数 $/$ 总样本数；
  \item 树最多 4 层、15 个决策节点；\textbf{全部 56 条残差样本仍参与期望缺口成本}，
        未把场景简化为少数均值曲线。
\end{enumerate}

\begin{figure}[H]
\centering
\begin{tikzpicture}[scale=1.0,every node/.style={font=\small},
   nd/.style={circle,draw,fill=blue!8,minimum size=6.5mm,inner sep=0pt},
   lf/.style={rectangle,draw,fill=green!10,minimum size=5mm,inner sep=2pt}]
  \node[nd] (r) at (0,0) {$0$};
  \node[nd] (a) at (2.2,1.3) {$1$};
  \node[nd] (b) at (2.2,-1.3) {$2$};
  \node[nd] (c) at (4.6,2.0) {$3$};
  \node[nd] (d) at (4.6,0.7) {$4$};
  \node[lf] (e) at (6.9,1.4) {$\vdots$};
  \node[lf] (f) at (6.9,-0.4) {$\vdots$};
  \draw (r)--(a) node[midway,above,sloped,font=\scriptsize]{$6$ 时修订};
  \draw (r)--(b);
  \draw (a)--(c) node[midway,above,sloped,font=\scriptsize]{$12$ 时修订};
  \draw (a)--(d);
  \draw (c)--(e) node[midway,above,sloped,font=\scriptsize]{$18$ 时};
  \draw (d)--(f);
  \node[font=\scriptsize,align=left] at (0,-2.3)
    {节点内\textbf{共享} $q,c,d,E$；\\ 仅 $e_{s,t}$ 按场景不同};
\end{tikzpicture}
\caption{经验场景树结构（示意）：顺着“新发布的预报修订”逐层分支，节点内共享控制变量}
\label{fig:q3-tree}
\end{figure}

\vspace{0.2em}
\noindent\textbf{为何“节点内共享变量”能保证非前视：}
因为控制变量 $q,c,d,E$ 是\textbf{同一个变量}跨越该节点内的所有场景，
物理上不可能针对“已经知道是哪个场景”来分别决策——
这在变量结构层面（而非仅通过额外约束）排除了前视。

\subsection{场景树随机规划模型}

以节点 $n$ 为单位，其概率为 $\pi_n$，控制变量为 $q_{n,t},c_{n,t},d_{n,t},E_{n,t}$：

\begin{equation}
\min\ J_3=\sum_tp_tg_t^0
+\sum_n\pi_n\sum_{t\in n}\bigl(1.5p_tu_{n,t}-0.5p_tv_{n,t}\bigr)
+\frac{1}{S}\sum_s\sum_t5p_te_{s,t}
+\underbrace{\text{次日前瞻成本}}_{\text{价值近似}}.
\label{eq:q3-obj}
\end{equation}

约束：

\begin{align}
&\text{（调整关系）}\quad q_{n,t}=g_t^0+u_{n,t}-v_{n,t},\qquad q,u,v\ge 0, \label{eq:q3-adj}\\
&\text{（节点需求）}\quad q_{n,t}+d_{n,t}-c_{n,t}+e_{s,t}\ \ge\ N_{s,t},\qquad e_{s,t}\ge0, \label{eq:q3-dem}\\
&\text{（计划下界）}\quad q_{n,t}+d_{n,t}-c_{n,t}\ \ge\ \hat N_t, \label{eq:q3-floor}\\
&\text{（SOC 与容量）}\quad E_{n,t}=E_{n,t-1}+0.9c_{n,t}-\frac{d_{n,t}}{0.9},\quad
   1200\le E_{n,t}\le10800, \label{eq:q3-soc}\\
&\text{（功率）}\quad 0\le c_{n,t},d_{n,t}\le\frac{5000}{6}, \label{eq:q3-pow}\\
&\text{（父子衔接）}\quad E_{n,\text{首}}=E_{\text{父},\text{末}}. \label{eq:q3-link}
\end{align}

\subsection{跨日价值与贪心执行}

\textbf{跨日前瞻}：每次优化至\textbf{次日 24 点}，把次日作为确定性价值近似，
避免“只顾当天便宜电价而提前耗空电池”。
不同未来节点使用其条件预报修订后的次日点预测。

\textbf{贪心执行}：只执行当前发布到下一允许更新前的控制量：

\begin{center}
\begin{tabular}{@{}c c@{}}
\toprule
决策时刻 & 本次执行区间 \\
\midrule
$0$:00  & $0$--$6$ 时 \\
$6$:00  & $6$--$12$ 时 \\
$12$:00 & $12$--$18$ 时 \\
$18$:00 & $18$--$24$ 时 \\
\bottomrule
\end{tabular}
\end{center}

\noindent\textbf{为何不把问题二的 $80\%$ 分位直接套用到本问：}
问题二的 $F^*=0.80$ 是在“储能固定、单期补救费为 $5$ 倍、无后续调整”
条件下导出的单期临界分位；而问题三中 0 点可从场景树\textbf{评估未来的调整机会}，
今天的“少买”可在 6/12/18 点以 $1.5$ 倍价补上，故有效风险水平随节点、时段而变，
不应硬编码单一分位。这正是不用报童模型而用场景树多阶段模型的原因。
（同时充放电的处理与问题一相同：由 \S\ref{sec:q1-tight} 的等 SOC 扰动自动消去，
程序对最终执行结果逐段检查互斥。）

\subsection{结果与边际分析}

\begin{table}[H]
\centering
\caption{问题三：更新频率对照与逐级边际收益（2 月 1 日--12 月 31 日）}
\label{tab:q3-abl}
\begin{tabular}{@{}l r r l@{}} 
\toprule
方案 & {总费用（元）} & {紧急电量（kWh）} & {增量节省（元）} \\
\midrule
方案 B 框架同口径复算 & 14389892.50 & 246099.33 & — \\
仅 0 点融合预报       & 14181445.03 & 218589.30 & 208447.47 \\
$0$、$6$ 点           & 14101371.47 & 186570.89 & 80073.56 \\
$0$、$6$、$12$ 点     & 14022358.10 & 175344.39 & 79013.36 \\
\textbf{$0$、$6$、$12$、$18$ 点（主方案）} & \bfseries 14022396.98 & 175369.95 & $-38.88$ \\
\bottomrule
\end{tabular}
\end{table}

\begin{table}[H]
\centering
\caption{问题三主方案费用分解}
\label{tab:q3-fee}
\begin{tabular}{@{}l r@{}}
\toprule
项目 & {金额（元）} \\
\midrule
原计划购电费      & 13345819.30 \\
增购费（$1.5$ 倍）  & 89947.28 \\
取消原费退款      & $-150170.13$ \\
取消违约费（$50\%$）& 75085.07 \\
净调整费（小计）    & 14862.22 \\
紧急购电费（$5$ 倍）& 661715.46 \\
\midrule
\textbf{总费用} & \bfseries 14022396.98 \\
\bottomrule
\end{tabular}
\end{table}

相对同口径方案 B，主方案\textbf{节省 $367\,495.52$ 元（$2.55\%$）}，
紧急购电量\textbf{减少 $28.74\%$}。
\textbf{结论}：$6$ 点与 $12$ 点各贡献约 $8$ 万元，收益随调整次数\textbf{显著边际递减}；
$18$ 点更新在本年度回测中\textbf{无可观测收益}（$-38.88$ 元），可视为可选更新。
该结论限于本数据与算法：$18$ 点仍可能通过次日预报影响跨日储能，
且小型经验树在有限样本下重建，不保证费用随信息增加单调下降。

\begin{figure}[htbp]
\centering
\includegraphics[width=0.88\textwidth]{04_q3_information_value.pdf}
\caption{固定电价情形下的逐级费用节省与紧急购电量。费用节省以本问同口径
历史预测对照为起点，依次引入 0 点融合预报和 6、12、18 点调整；
右图按相同顺序比较紧急购电量。}
\label{fig:q3-info-value}
\end{figure}

图 \ref{fig:q3-info-value} 把“信息更新的价值”与“更新次数的增加”分开呈现：
仅引入 0 点融合预报即节省 $208\,447.47$ 元；
随后增加 6 点与 12 点更新各节省约 $8$ 万元；增加 18 点后费用反而上升 $38.88$ 元。
这说明\textbf{预报融合本身的价值最大，而更新次数存在明显边际递减}。

\begin{figure}[htbp]
\centering
\includegraphics[width=0.88\textwidth]{05_q3_forecast_fusion.pdf}
\caption{发布后六小时光伏预测误差及融合权重。(a)(b) 分别比较日间和夜间的
历史预测、发布预报及在线融合 MAE，两面板纵轴尺度不同；
(c) 展示各发布时间对应的发布预报权重。}
\label{fig:q3-fusion}
\end{figure}

6--12 时段的 MAE 由历史预测的 $309.24$ kW 降至融合预测的 $169.41$ kW，
12--18 时段由 $288.99$ kW 降至 $133.68$ kW。夜间绝对误差较小，单独绘制以保留其变化尺度。
融合权重随历史已成熟样本更新，并包含偏差校正，因此融合输出不只是当期两条预报的
固定加权平均。\textbf{MAE 改善本身不等同于等比例的费用下降}，
经济效果须结合图 \ref{fig:q3-info-value} 的更新频率对照判断。

\begin{figure}[htbp]
\centering
\includegraphics[width=0.88\textwidth]{06_q3_rolling_adjustment.pdf}
\caption{四个指定日的日内计划修正与储能状态。左列为最终执行购电相对 0 点原计划的
功率差，正值为增购、负值为减购；右列为储电量，底部橙色短线标记紧急购电
大于数值容差的时段。竖虚线表示 6、12、18 点发布时刻。}
\label{fig:q3-adjust}
\end{figure}

日内更新后的购电修正相对原计划较小，直接叠加两条绝对购电曲线容易掩盖变化；
以原计划为零基准后，可清楚识别各执行区间内的增购与减购。
储电量按跨日滚动过程衔接，指定日的首末电量不要求逐日相等。
紧急购电标记只表示发生时间，不表示缺口大小，
也不能据其与储能状态的相邻关系断言单一原因。

\FloatBarrier

\subsection{敏感性分析与性能边界}
\label{sec:q3-sens}

\textbf{完美预见下界}：假设每天 $0$ 点即精确知道当天实际净负载（消除全部预报误差），
在同一物理约束与计费规则下求解得下界 $1\,223.16$ 万元（紧急购电费为 $0$），
与主方案 $1\,402.24$ 万元相差 $179.08$ 万元——\textbf{该缺口全部来自预报误差，而非建模策略}。

\textbf{融合权重的饱和性}：逐日自适应融合的加权 MAE 为 93.68 kW，
\textbf{优于含未来信息泄漏的静态最优线性组合}（95.08 kW），
说明权重估计已实质达到该模型族的精度上界。

\textbf{参数扫描（30 余组）}：对场景缩放、残差窗口、收缩参数、计划下界、多源信息等
做同年度扫描，仅“场景缩放 $k=0.92$”有正收益（约 $0.47$ 万元，$0.03\%$），
其余全部无效或变差（窗口 $W=28$ 时 $+5.31$ 万元、计划下界 $0.5\times$ 时 $+12.15$ 万元）。
未把 $k=0.92$ 写入正式结果——同年度扫描存在过拟合风险，收益量级可忽略。

\textbf{注意到}：（1）\textbf{弃置电量并非浪费}——全年弃置 $267.66$ 万 kWh，
但放松计划下界反而使总费用上升，原因是“多买的电”是“低价购电 $+$ 储能套利”的必要冗余；
（2）\textbf{增加信息源可能有害}——“上一期发布”虽合法可用，但纳入融合后费用上升
$+0.66$ 万元，因多次发布高度相关，额外信息源主要引入估计噪声。

\noindent\textbf{结论的适用范围}：以上只能说明“在已测试模型族与同年度回测下未发现更优”，
\textbf{不足以宣称全局最优}；引入外部气象数据仍有提升余地。
